% Instructor's original relativistic notes pp. 6--8: Pauli--Dirac representation and limits.
\section{The Nonrelativistic Limit of the Dirac Equation}
\subsection*{1) Free-particle energy and the rest-energy phase}
For a free particle, the eigenvalues of the Dirac Hamiltonian satisfy
\begin{equation}
 E_\pm=\pm\sqrt{m^2c^4+c^2\vect p^{\,2}}.
\end{equation}
Near the positive-energy rest branch, $|\vect p|\ll mc$,
\begin{equation}
\begin{aligned}
E_+&=mc^2\sqrt{1+\frac{\vect p^{\,2}}{m^2c^2}}\\
&=mc^2+\frac{\vect p^{\,2}}{2m}-\frac{\vect p^{\,4}}{8m^3c^2}
  +O\!\left(\frac{p^6}{m^5c^4}\right).
\end{aligned}
\end{equation}
The constant rest energy creates a rapid overall phase. Remove it by writing
\begin{equation}
 \psi(\vect r,t)=e^{-\ii mc^2t/\hbar}
 \begin{pmatrix}\varphi(\vect r,t)\\\chi(\vect r,t)\end{pmatrix}.
\end{equation}
In the Pauli--Dirac representation, $\op H_{\mathrm D}$ has block form
\begin{equation}
 \op H_{\mathrm D}=\begin{pmatrix}mc^2\mathbf1_2&c\vect\sigma\cdot\op{\vect p}\\
 c\vect\sigma\cdot\op{\vect p}&-mc^2\mathbf1_2\end{pmatrix}.
\end{equation}
Substitution into the Dirac equation gives two coupled Pauli-spinor equations:
\begin{equation}\label{eq:dirac-two-spinors}
\boxed{\begin{aligned}
\ii\hbar\partial_t\varphi&=c\vect\sigma\cdot\op{\vect p}\,\chi,\\
\ii\hbar\partial_t\chi&=c\vect\sigma\cdot\op{\vect p}\,\varphi-2mc^2\chi.
\end{aligned}}
\end{equation}
The large rest-energy separation makes the lower spinor small on the positive-energy branch. When kinetic energies and the time variation of $\chi$ are small compared with $2mc^2$, the second equation gives
\begin{equation}
 \chi\simeq\frac{\vect\sigma\cdot\op{\vect p}}{2mc}\varphi.
\end{equation}
Inserting this into the first equation produces
\begin{equation}
 \ii\hbar\partial_t\varphi
 \simeq\frac1{2m}(\vect\sigma\cdot\op{\vect p})^2\varphi
 =\frac{\op{\vect p}^{\,2}}{2m}\varphi.
\end{equation}
Thus the positive-energy Dirac equation reduces to the free two-component Pauli equation. The lower component has a controlled suppression of order $|\vect p|/(mc)$. The approximation selects positive-energy, low-momentum states; the complete relativistic theory still contains antiparticle solutions.

\subsection*{2) Recovering spin coupling for a charged particle}
Under minimal electromagnetic coupling, replace $\op{\vect p}$ by
$\op{\vect\pi}=\op{\vect p}-(q/c)\vect A$ and add $q\phi$. After extracting the rest-energy phase, the equations become
\begin{equation}
\begin{aligned}
\ii\hbar\partial_t\varphi&=q\phi\,\varphi
       +c\vect\sigma\cdot\op{\vect\pi}\,\chi,\\
\ii\hbar\partial_t\chi&=c\vect\sigma\cdot\op{\vect\pi}\,\varphi
      +(q\phi-2mc^2)\chi.
\end{aligned}
\end{equation}
At leading nonrelativistic order,
\begin{equation}
 \chi\simeq\frac{\vect\sigma\cdot\op{\vect\pi}}{2mc}\varphi,
\end{equation}
so the upper equation reads
\begin{equation}\label{eq:pauli-from-dirac}
\boxed{\ii\hbar\partial_t\varphi=
\left[\frac1{2m}(\vect\sigma\cdot\op{\vect\pi})^2+q\phi\right]\varphi
+\text{higher-order relativistic corrections}.}
\end{equation}
Using the commutator of the kinetic momentum components obtained in the first section yields
\begin{equation}
\ii\hbar\partial_t\varphi=\left[
\frac{\op{\vect\pi}^{\,2}}{2m}
-\frac{q\hbar}{2mc}\vect\sigma\cdot\vect B+q\phi\right]\varphi
+\cdots.
\end{equation}
This is precisely the Pauli Hamiltonian in Eq.~\eqref{eq:pauli-expanded}. In particular, the spin magnetic interaction follows from the matrix factorization of relativistic energy. Higher-order terms bring additional effects such as relativistic kinetic corrections and, for inhomogeneous fields, spin--orbit and Darwin terms.

\section{The Relation among the Schr\"odinger, Pauli, and Dirac Equations}
The preceding derivations connect three familiar descriptions. The scalar Schr\"odinger equation is the nonrelativistic theory of orbital motion. Enlarging its state space to a Pauli spinor incorporates spin $1/2$ and magnetic coupling. The relativistic Dirac equation describes a four-component spinor; at small velocities and after factoring out the rest energy, its positive-energy sector returns the Pauli equation. The Klein--Gordon equation results directly from the relativistic energy relation, and for scalar particles it admits a quantum-field interpretation.

\begin{center}
\begin{tikzpicture}[>=Latex,thick,
 node distance=17mm and 28mm,
 every node/.style={font=\small,align=center}]
 \node (sch) [draw,rounded corners,minimum width=36mm,minimum height=11mm]
    {Schr\"odinger\\scalar wavefunction};
 \node (pau) [draw,rounded corners,right=of sch,minimum width=36mm,minimum height=11mm]
    {Pauli\\two-component spinor};
 \node (kg) [draw,rounded corners,below=of sch,minimum width=36mm,minimum height=11mm]
    {Klein--Gordon\\scalar relativistic field};
 \node (dir) [draw,rounded corners,below=of pau,minimum width=36mm,minimum height=11mm]
    {Dirac\\four-component spinor};
 \draw[->] (sch.east)--node[above,font=\scriptsize] {include spin}(pau.west);
 \draw[->] (sch.south)--node[left,font=\scriptsize,align=right] {relativistic\\dispersion}(kg.north);
 % Distinct side paths for the two different operations; neither crosses a box or another label.
 \draw[->] (pau.south) to[out=-105,in=105] node[below left,font=\scriptsize,pos=.52,align=right]
      {relativistic\\extension} (dir.north);
 \draw[->] (dir.north) to[out=75,in=-75] node[right,font=\scriptsize,pos=.5,align=left]
      {low-energy\\limit} (pau.south);
\end{tikzpicture}
\end{center}
In particular, the two-component Pauli equation can be viewed both as the spinful extension of nonrelativistic Schr\"odinger theory and as a systematically controlled limit of the Dirac equation.